\section{Benchmark Evidence: por qué elegir tabular data y cómo leer las tablas}

Si solo se prueba en benchmarks que las redes profundas dominan con facilidad, es difícil saber si la dataset-level attention aporta valor real. Los autores eligen tabular data porque abarca clasificación y regresión, campos continuos y discretos, y desde cientos hasta decenas de millones de muestras, además de que el tree-based boosting ha sido fuerte durante mucho tiempo; esto obliga a comparar NPT con baselines no neuronales que de verdad compiten.

\lecturefigure{slide-22-tabular-domain.jpg}{Tabular benchmark: la escala, los tipos de campo, las tareas y los baselines fuertes son muy diversos.}{Video oficial de Stanford Online 00:37:16--00:38:07.}

\begin{knowledgebox}{Referencia rápida de baselines y métricas}
\begin{tabularx}{\textwidth}{L{0.19\textwidth}L{0.25\textwidth}>{\raggedright\arraybackslash}X}
\toprule
Nombre & Mecanismo & Por qué es necesario compararlo \\
\midrule
XGBoost / LightGBM & gradient-boosted decision trees eficientes & Baseline fuerte de la industria para tabular data; maneja bien las divisiones no lineales y las características faltantes o dispersas. \\
CatBoost & ordered boosting orientado a categorical features & Evita que NPT obtenga una ventaja falsa solo por tratar mejor los campos categóricos. \\
Random Forest / GBDT & Modelos de árboles con bagging o boosting & Cubre la ensemble family clásica. \\
MLP & parametric neural baseline que procesa cada fila de forma independiente & Verifica si la mejora proviene del mecanismo entre datapoints. \\
TabNet & attentive neural model orientado a tabular data & Compara con una domain-specific neural architecture. \\
k-NN & local lookup con distancia fija & Contrasta si el «learned non-parametric lookup» supera a los vecinos simples. \\
AUROC / Accuracy / RMSE & Ordenamiento en clasificación binaria, tasa de aciertos en clasificación multiclase, error de regresión & Las métricas de los tres tipos de tarea tienen direcciones distintas; el artículo las unifica convirtiéndolas en method rank. \\
\bottomrule
\end{tabularx}
\end{knowledgebox}

\teachervoice{Los autores no eligieron el tabular domain solo porque los datos fueran convenientes, sino porque las general-purpose deep nets suelen perder aquí frente al boosting; si NPT solo superara a un MLP débil, eso no bastaría para sostener su afirmación de generalidad.}

\teachervoice{Al equipo le sorprendió la robustness de NPT en small datasets, y eso sin realizar una cantidad extrema de ajustes de hyperparameters. Vale la pena registrar esta observación verbal, pero no sustituye una comparación rigurosa del tuning-budget.}

\subsection{Average rank no es pooled accuracy}

La página de resultados agrupa los diez UCI datasets por tarea: cuatro de binary classification usan AUROC, dos de multi-class classification usan accuracy y cuatro de regression usan RMSE. Para cada dataset primero se ordenan los methods y luego se promedia dentro del mismo grupo de tareas; por ello, un valor menor es mejor, y no se pueden comparar directamente columnas distintas como si fueran la misma magnitud física.

\lecturefigure{slide-23-tabular-benchmarks.jpg}{Rango promedio en el UCI benchmark y primeros resultados en CIFAR-10.}{Video oficial de Stanford Online 00:38:08--00:40:44.}

\begin{table}[H]
\centering
\small
\caption{Rangos promedio clave de la Table 1 del artículo; un valor menor es mejor.}
\begin{tabularx}{\textwidth}{L{0.22\textwidth}L{0.23\textwidth}L{0.23\textwidth}X}
\toprule
Método & Binary AUROC rank & Multi-class accuracy rank & Regression RMSE rank \\
\midrule
NPT & $2.50\pm0.87$ & $2.50\pm0.50$ & $3.25\pm1.31$ \\
CatBoost & $2.75\pm0.85$ & $3.50\pm0.50$ & $3.00\pm0.91$ \\
XGBoost & $4.75\pm1.25$ & $2.50\pm1.50$ & $3.25\pm0.63$ \\
MLP & $5.75\pm1.49$ & $3.00\pm2.00$ & $5.00\pm1.22$ \\
k-NN & $8.25\pm0.48$ & $8.50\pm0.50$ & $8.75\pm0.25$ \\
\bottomrule
\end{tabularx}
\end{table}

\begin{knowledgebox}{Lectura de la figura: primero el intervalo competitivo, luego el número de primeros lugares}
En binary y multi-class, el rank promedio de NPT está en el grupo de cabeza, y en regression queda cerca de CatBoost/XGBoost; el artículo también informa que NPT es top performer en cuatro de los diez tabular datasets. La conclusión más prudente es «competitivo en varios tipos de tarea», no «supera de manera uniforme a todos los tree methods».
\end{knowledgebox}

\begin{warningbox}{Qué oculta el rango promedio}
Diez datasets son pocos, y cada grupo de tareas tiene solo de dos a cuatro conjuntos de datos; el rank pierde el performance gap absoluto y además depende del tuning protocol y de la dataset selection. El error estándar proviene de la rank variation entre datasets; no es el intervalo de confianza de una sola accuracy medida sobre una muestra grande.
\end{warningbox}

\teachervoice{En la sesión de Q\&A se dedicó bastante tiempo a aclarar «4 of 10», «2 of 10» y el standard error. Por eso las notas deben dejar claro que estas cifras describen primero la agrupación por dataset/task y solo después el rango promedio de los methods; a partir de los decimales en apariencia pequeños no puede inferirse una ventaja significativa general.}

\subsection{Los límites correctos del image result}

El artículo también prueba NPT en CIFAR-10 y MNIST. El resultado en CIFAR-10 que se muestra en la clase usa un CNN encoder más NPT y reporta una accuracy de 93.7\%; esto indica que la datapoint interaction entre distintas images puede integrarse en un pipeline de visión, pero no demuestra que NPT sustituya al backbone visual. Los resultados con linear patching del apéndice del artículo son más débiles, lo que también refleja que el input embedding sigue influyendo en el desempeño.

\begin{warningbox}{«Formato de entrada general» no equivale a «el front-end es irrelevante»}
Que cualquier modality pueda hacer reshape a una matrix solo indica que la interfaz puede recibirla; si se conserva la estructura local, si existe un encoder adecuado y cómo se configuran el volumen de datos y la augmentation siguen determinando la optimization y la sample efficiency. La contribución de NPT es agregar un mecanismo entre datapoints, no eliminar todo el conocimiento propio de cada modality.
\end{warningbox}

\subsection{Resumen de la sección}

NPT entra en la zona competitiva de los baselines fuertes en diversas tabular tasks y muestra la viabilidad de la image extension. La evidencia respalda que es «competitivo», pero el rank promedio, el número limitado de datasets y los distintos front-ends no nos permiten afirmar que sea el mejor en todos los escenarios.
